\subsection{Two further reductions for the four-object seventeen-arrow scope}

\begin{lemma}[Proved: a height-one support is a retract]
\label{CP2:A-height-one-support}
Let $C$ have acyclic support with no composable pair of cross-object
arrows. Its support poset $P$ is a categorical retract of $C$.
Consequently, if the undirected support graph has positive cycle rank,
then $H_1(BC;\mathbb Z)\ne0$.
\end{lemma}
\begin{proof}
The support functor $C\to P$ sends every corner to its identity and
every cross arrow to the corresponding unique poset arrow. Choose
one representative in each nonempty cross hom-set. These choices,
together with the corner identities, define a functor $P\to C$:
there are no composable pairs of nonidentity arrows in $P$ whose
composition would impose an additional condition. Its composite
with $C\to P$ is the identity. The nerve of $P$ is exactly its
undirected support graph. The induced retraction on first homology
proves the claim. In particular every four-object height-one crown
is excluded, regardless of its corner monoids or hom-set sizes.
\end{proof}

\begin{proposition}[Proved: the eight remaining large-corner profiles]
\label{CP2:A-four17-large-corner-profiles}
No four-object category of any of the following hom-cardinality
matrices has type $\CP^2$:
\[
\begin{gathered}
\begin{pmatrix}
1&0&0&1\\0&a&s&t\\0&0&b&0\\0&0&0&c
\end{pmatrix},
\quad 8\leq a\leq9,\quad b,c\leq2,\quad s,t\geq2,
\quad a+b+c+s+t+2=17;\\[3pt]
\begin{pmatrix}
1&0&0&2\\0&1&s&2\\0&0&b&0\\0&0&0&8
\end{pmatrix},
\quad (b,s)=(2,1)\text{ or }(1,2);\\[3pt]
\begin{pmatrix}
1&0&1&1\\0&8&0&3\\0&0&1&1\\0&0&0&1
\end{pmatrix}.
\end{gathered}
\]
There are respectively five, two, and one matrices in these families.
\end{proposition}
\begin{proof}
We use actual unions of full subcategory nerves. If every simplex
lies in $ND$ or $NE$, their intersection is a full subcategory $F$,
and $H_j(BF;\mathbb Q)=0$ for $j>0$, Mayer--Vietoris gives
\[
 H^j(BC;\mathbb Q)\cong
 H^j(BD;\mathbb Q)\oplus H^j(BE;\mathbb Q),\qquad j\geq2.
\]
Thus bounded rational homology of $BC$ implies the same for the two
pieces in degrees at least two. Moreover, products vanish on $BC$
if they vanish on both pieces, because restriction is injective
in degree four. All the intersections used below have corner
order at most two and hence are rationally acyclic.

For the first family take $D$ on objects $1,2,3$ and $E$ on $0,3$.
Every simplex belongs to one of these pieces and $F=\{3\}$.
The edge $E$ has a singleton cross hom-set and a trivial source,
so $BE$ is contractible. The category $D$ is a source star with
center order $a\leq9$ and leaves of orders at most two. Exact
orbit comparison at the leaves, Lemma~\ref{CP2:A-high-orbits},
makes the leaves trivial by a rational cohomology-ring
equivalence. Identifying these two trivial leaves gives an
ordinary act cone over the center, whose space is the reduced
star wedged with a circle, as in
Lemma~\ref{CP2:A-star-leaf-reductions}(3). It has bounded rational
homology, so Theorem~\ref{CP2:A-bounded-nine-acts} kills all its
degree-two products. Those of $BD$, and then of $BC$, vanish.

For the second family use $D$ on $0,1,3$ and $E$ on $1,2$,
intersecting at the trivial corner $1$. The first piece is the
opposite of the same bounded-nine act construction: its center
has order eight and its two leaves are trivial. The second piece
is either a singleton edge or a two-arrow edge with trivial
corners, so has zero second rational cohomology. The same
restriction argument applies.

For the third matrix use $D$ on $0,2,3$ and $E$ on $1,3$.
The first category is the three-element chain and is contractible;
the intersection is the trivial corner $3$. The second category
is an ordinary act cone over an eight-element monoid. Boundedness
and the bounded-nine act theorem again give zero products.
In every case this contradicts the nonzero rational square on
$\CP^2$.
\end{proof}
